\documentclass[11pt,a4paper]{article}
\usepackage[utf8]{inputenc}
\usepackage[T1]{fontenc}
\usepackage[french]{babel}
\usepackage{amsmath}
\usepackage{amssymb}
\usepackage{listings}
\usepackage{xcolor}
\usepackage{graphicx}
\usepackage{geometry}
\geometry{margin=2.5cm}

% Configuration des blocs de code
\lstset{
    basicstyle=\ttfamily\footnotesize,
    keywordstyle=\color{blue}\bfseries,
    commentstyle=\color{green!60!black},
    stringstyle=\color{red},
    showstringspaces=false,
    frame=single,
    numbers=left,
    numberstyle=\tiny\color{gray},
    breaklines=true,
    tabsize=4
}

\title{TD3 Correction \\ \large SN360/SN361 : Introduction à la conception des circuits numériques}
\author{Grenoble INP - Esisar - EI3-2024/2025}
\date{}

\begin{document}

\maketitle

\section*{Exercice 1 : Conception d'un compteur électronique}

En électronique, un compteur est un composant essentiel[cite: 582]. Considérons ici un compteur 8 bits comportant[cite: 582]:
\begin{itemize}
    \item Une entrée d'activation synchrone appelée \texttt{Enable}[cite: 583].
    \item Un signal de reset asynchrone nommé \texttt{rst}[cite: 584].
    \item Une incrémentation qui s'effectue sur les fronts montants de l'horloge \texttt{clk}[cite: 585].
\end{itemize}

\subsection*{Questions :}

\paragraph{1) Combien de bascules sont nécessaires pour implémenter ce compteur ?}
8 bascules, chaque bit étant représenté par une bascule[cite: 587].

\paragraph{2) Dessinez un chronogramme illustrant le fonctionnement du compteur.}
Il s'agit d'un exemple de simulation avec Vivado pour un compteur utilisant une horloge d'une période de 10 ns[cite: 588]. *(Le chronogramme affiche l'évolution de \texttt{count} de 00 à 13 en hexadécimal à chaque front d'horloge lorsque \texttt{Enable} est actif [cite: 589, 592-602, 632-640]).*

\paragraph{3) Écrivez le code de ce compteur en VHDL et en Verilog.}

\textbf{VHDL :}
\begin{lstlisting}[language=VHDL]
library IEEE;
use IEEE.STD_LOGIC_1164.ALL;
use IEEE.STD_LOGIC_UNSIGNED.ALL;
use IEEE.NUMERIC_STD.ALL;

entity compteur is
Port (
    clk    : in STD_LOGIC;  -- horloge
    rst    : in STD_LOGIC;  -- reset
    Enable : in STD_LOGIC;  -- activation
    count  : out STD_LOGIC_VECTOR (7 downto 0) -- sortie sur 8 bits
);
end compteur;

architecture Behavioral of compteur is
    signal count_internal: unsigned(7 downto 0) := (others => '0');
begin
    process (clk, rst)
    begin
        if rst = '1' then
            count_internal <= (others => '0');
        elsif rising_edge(clk) then
            if Enable = '1' then
                count_internal <= count_internal + 1;
            end if;
        end if;
    end process;
    count <= std_logic_vector(count_internal);
end Behavioral;
\end{lstlisting}

\textbf{Verilog :}
\begin{lstlisting}[language=Verilog]
module compteur (
    input wire clk,
    input wire rst,
    input wire Enable,
    output reg [7:0] Count 
);

always @(posedge clk or posedge rst) begin
    if (rst) begin
        Count <= 8'b0;
    end else if (Enable) begin
        Count <= Count + 1;
    end
end
endmodule
\end{lstlisting}

\paragraph{4) Créer le TestBench pour le compteur en VHDL et Verilog :}
Le TestBench permet de vérifier le bon fonctionnement du système[cite: 660]. Dans l'exemple suivant, on teste la réaction du système à une horloge avec une période de 10 ns[cite: 661]. Le signal reset et le signal enable sont respectivement à 1 et 0 pendant 50 ns[cite: 662]. Ensuite, reset passe à 0 et enable à 1 pendant 200 ns[cite: 663]. Enfin, les deux signaux sont positionnés à 1[cite: 664].

\textbf{TestBench VHDL :}
\begin{lstlisting}[language=VHDL]
library IEEE;
use IEEE.STD_LOGIC_1164.ALL;
use IEEE.STD_LOGIC_UNSIGNED.ALL;

entity compteur_tb is
end compteur_tb;

architecture Behavioral of compteur_tb is
    component compteur
    Port (
        clk    : in STD_LOGIC;
        rst    : in STD_LOGIC;
        Enable : in STD_LOGIC;
        count  : out STD_LOGIC_VECTOR (7 downto 0)
    );
    end component;

    signal clk    : STD_LOGIC := '0';
    signal rst    : STD_LOGIC := '0';
    signal Enable : STD_LOGIC := '0';
    signal count  : STD_LOGIC_VECTOR (7 downto 0);
begin
    uut: compteur Port map (
        clk    => clk,
        rst    => rst,
        Enable => Enable,
        count  => count
    );

    clk_process: process
    begin
        clk <= not(clk);
        wait for 5 ns;
    end process;

    stim_proc: process
    begin
        rst <= '1';
        Enable <= '0';
        wait for 50 ns;
        rst <= '0';
        Enable <= '1';
        wait for 200 ns;
        Enable <= '1';
        rst <= '1';
        wait;
    end process;
end Behavioral;
\end{lstlisting}

\textbf{TestBench Verilog :}
\begin{lstlisting}[language=Verilog]
`timescale 1ns/1ps

module tb_compteur();
    reg clk = 1;
    reg rst = 0;
    reg Enable = 0;
    wire [7:0] Count;

    compteur uut (
        .clk(clk),
        .rst(rst),
        .Enable(Enable),
        .Count(Count)
    );

    initial begin
        forever #5 clk = ~clk;
    end

    initial begin
        rst = 1;
        Enable = 0;
        #50 rst = 0;
        Enable = 1;
        #200 rst = 1;
        Enable = 1;
    end
endmodule
\end{lstlisting}

\paragraph{Amélioration :}
Nous souhaitons maintenant étendre le fonctionnement du compteur avec deux nouvelles entrées[cite: 739]:
\begin{itemize}
    \item \texttt{up} : Si active, le compteur s'incrémente[cite: 740].
    \item \texttt{down} : Si active, le compteur se décrémente[cite: 741].
\end{itemize}

\textbf{Modifiez le code VHDL pour intégrer cette nouvelle fonctionnalité :}
\begin{lstlisting}[language=VHDL]
library IEEE;
use IEEE.STD_LOGIC_1164.ALL;
use IEEE.NUMERIC_STD.ALL;

entity compteur_ext is
Port (
    clk    : in STD_LOGIC;
    rst    : in STD_LOGIC;
    Enable : in STD_LOGIC;
    up     : in STD_LOGIC;
    down   : in STD_LOGIC;
    count  : out STD_LOGIC_VECTOR (7 downto 0)
);
end compteur_ext;

architecture Behavioral of compteur_ext is
    signal count_internal: unsigned (7 downto 0) := (others => '0');
begin
    process (clk, rst)
    begin
        if rst = '1' then
            count_internal <= (others => '0'); -- Reinitialisation du compteur
        elsif rising_edge(clk) then
            if Enable = '1' then
                if up = '1' and down = '0' then
                    count_internal <= count_internal + 1; -- Incrementation
                elsif down = '1' and up = '0' then
                    count_internal <= count_internal - 1; -- Decrementation
                else 
                    count_internal <= (others => '0'); -- si up et down actifs en meme temps
                end if;
            end if;
        end if;
    end process;
    count <= std_logic_vector(count_internal);
end Behavioral;
\end{lstlisting}

\vspace{1cm}

\section*{Exercice 2 : Conception d'un registre à décalage SIPO de N bits}

\paragraph{1) On dispose de la description VHDL d'une bascule D (flip-flop) dont l'entité est donnée ci-dessous[cite: 805]. Ecrivez le testbench de ce composant[cite: 806].}

\begin{lstlisting}[language=VHDL]
entity bascule is
port(d, clk, reset, enable: in std_logic;
     q: out std_logic);
end entity bascule;
\end{lstlisting}

Voici un TestBench conçu pour tester le système. L'objectif est de vérifier un maximum d'entrées, voire toutes si possible, afin d'observer la réaction de la sortie et de s'assurer qu'elle correspond aux attentes[cite: 811].

\begin{lstlisting}[language=VHDL]
library IEEE;
use IEEE.STD_LOGIC_1164.ALL;

entity tb_bascule is
end tb_bascule;

architecture bhv of tb_bascule is
    signal tb_clk    : std_logic := '0';
    signal tb_d      : std_logic := '0';
    signal tb_enable : std_logic := '0';
    signal tb_reset  : std_logic := '0';
    signal tb_q      : std_logic;

    component bascule
    port (
        d, clk, reset, enable: in std_logic;
        q: out std_logic
    );
    end component;

begin
    i1: bascule port map (
        d      => tb_d,
        clk    => tb_clk,
        reset  => tb_reset,
        enable => tb_enable,
        q      => tb_q
    );

    -- Il est possible de tester le signal d'horloge (clk) sans utiliser de processus.
    tb_clk <= not tb_clk after 5 ns; 

    stim_proc: process
    begin
        tb_reset <= '1';
        tb_enable <= '0';
        tb_d <= '0';
        wait for 10 ns;
        
        tb_reset <= '0';
        wait for 10 ns;
        
        tb_d <= '1';
        wait for 20 ns;
        
        tb_enable <= '1';
        wait for 10 ns;
        
        tb_d <= '0';
        wait for 20 ns;
        
        tb_reset <= '1';
        wait for 10 ns;
        
        tb_reset <= '0';
        wait for 20 ns;
        
        wait;
    end process;
end bhv;
\end{lstlisting}

\paragraph{2) On veut construire un registre à décalage de 2 bits SIPO (Serial In - Parallel Out) en utilisant uniquement 2 bascules D[cite: 853].}
Ce registre à décalage comporte[cite: 854]:
\begin{itemize}
    \item un simple bit d'entrée (\texttt{din}) [cite: 855]
    \item une entrée d'autorisation de décalage (\texttt{enable}) active à l'état haut [cite: 856]
    \item une sortie parallèle (\texttt{dout}) sur 2 bits [cite: 857]
    \item une remise à zéro (\texttt{reset}) synchrone active à l'état haut[cite: 858].
\end{itemize}

\textbf{2.1. Dessinez le schéma de ce registre}
*(Le schéma est à inclure ici à l'aide du fichier SVG généré).*

\textbf{2.2. Ecrivez la description structurelle du registre SIPO (entité + architecture) [cite: 867]}
\begin{lstlisting}[language=VHDL]
library IEEE;
use IEEE.std_logic_1164.all;

entity SIPO is
port( 
    din, enable, clk, reset: in std_logic;
    dout: out std_logic_vector(0 to 1)
);
end SIPO;

architecture struct of SIPO is
    component bascule
    port(
        d, clk, reset, enable: in std_logic;
        q: out std_logic
    );
    end component;

    signal dout_temp: std_logic_vector(0 to 1);
begin
    i1: bascule port map(
        d => din, enable => enable, clk => clk, reset => reset, q => dout_temp(0)
    );
    
    i2: bascule port map (
        d => dout_temp(0), enable => enable, clk => clk, reset => reset, q => dout_temp(1)
    );
    
    dout <= dout_temp;
end struct;
\end{lstlisting}

\paragraph{3) Pour que le registre soit facilement réutilisable dans différents projets, on souhaite le rendre plus générique[cite: 889].} 
On veut qu'il soit paramétrable pour supporter différentes tailles de vecteurs de sortie[cite: 890]. On veut donc construire un registre à décalage SIPO de N bits[cite: 891]. Ecrivez l'entité de ce registre[cite: 892].

\begin{lstlisting}[language=VHDL]
library IEEE;
use IEEE.std_logic_1164.all;

entity SIPO_n is
generic (N: integer := 5);
port (
    din    : in std_logic;
    enable : in std_logic;
    clk    : in std_logic;
    reset  : in std_logic;
    dout   : out std_logic_vector(0 to N-1)
);
end SIPO_n;

architecture struct of SIPO_n is
    component bascule
    port (
        d      : in std_logic;
        clk    : in std_logic;
        reset  : in std_logic;
        enable : in std_logic;
        q      : out std_logic
    );
    end component;

    signal dout_temp: std_logic_vector(0 to N);
begin
    bascules_gen: for I in 0 to N-1 generate
    begin
        bascule_inst: bascule port map (
            d      => dout_temp(I),
            enable => enable,
            clk    => clk,
            reset  => reset,
            q      => dout_temp(I+1)
        );
    end generate;
    
    dout_temp(0) <= din;
    dout <= dout_temp(1 to N);
end architecture;
\end{lstlisting}

\textbf{TestBench pour SIPO\_n[cite: 941]:}
\begin{lstlisting}[language=VHDL]
library IEEE;
use IEEE.std_logic_1164.all;

entity tb_SIPO_n is
end tb_SIPO_n;

architecture bhv of tb_SIPO_n is
    signal tb_clk    : std_logic := '1';
    signal tb_din, tb_enable, tb_reset: std_logic;
    signal tb_dout   : std_logic_vector(0 to 4);

    component SIPO_n
    generic(N: integer := 5);
    port(
        din, enable, clk, reset: in std_logic;
        dout: out std_logic_vector(0 to N-1)
    );
    end component;
begin
    i1: SIPO_n 
    generic map( N => 5 )
    port map (
        din    => tb_din,
        enable => tb_enable, 
        clk    => tb_clk,
        reset  => tb_reset,
        dout   => tb_dout
    );

    tb_clk <= not(tb_clk) after 5 ns;
    
    tb_din <= '0', '1' after 5 ns, '0' after 25 ns, '1' after 45 ns;
    tb_enable <= '0', '1' after 15 ns, '0' after 25 ns, '1' after 35 ns;
    tb_reset <= '0', '1' after 55 ns, '0' after 65 ns;
end bhv;
\end{lstlisting}

\end{document}